

\section{Sixteenth Sunday after Trinity: A Faithful Friend in Distress and Death}


\begin{quote}

John 11:32-45. When Mary now came where Jesus was, and saw him, she fell down at his feet and said to him, Lord, if you had been here, my brother would not have died. When Jesus now saw her weeping, and the Jews also weeping who had come with her, he was deeply moved in spirit and shaken, and said, Where have you laid him? They said to him, Lord, come and see. Jesus wept. Then said the Jews, See how he loved him! But some of them said, Could not he who opened the eyes of the blind man also have caused this man not to die? Then Jesus, again deeply moved, came to the tomb. Now it was a cave, and a stone lay upon it. Jesus said, Take away the stone. Martha, the sister of the dead man, says to him, Lord, he already stinks, for he has lain there four days. Jesus said to her, Did I not tell you that if you could believe, you would see the glory of God? Then they took away the stone where the dead man had been laid. But Jesus lifted up his eyes and said, Father, I thank you that you have heard me. Yet I knew that you always hear me; but for the sake of the people who stand around I said it, that they might believe that you have sent me. And when he had said this, he cried with a loud voice, Lazarus, come out here! And the dead man came out, bound hand and foot with grave-clothes, and his face was wrapped with a cloth. Jesus said to them, Loose him, and let him go. Then many of the Jews who had come to Mary, and had seen what Jesus did, believed in him.

\end{quote}

\bigskip

Four things are required if a friend is to be a true, perfect friend, one in whom I may securely place the whole trust of my heart. And who would not wish to possess such a one? The first is that he can truly feel with me, that he can wholly and fully understand me in all the innermost need and want of my heart, in sorrow and in joy, in distress and in death. The second is that he loves me so highly that he is ready to give up everything, were it even his life, in order to help me and make me glad and blessed. The third is that he truly also has power to do what his love and faithful friendship bid him do, that, in other words, he not only gladly will, but also can. The fourth is that such friendship has its ground in that which is eternal and unchangeable with God.

Where among men is such a friend to be found? Jesus is such a friend, and reveals himself thus in our text.

For in the consciousness of many, Jesus has receded so far away into the depth of God’s power and eternity that for them he stands almost only as a doctrine of atonement. His personality vanishes into what is purely dogmatic. That may be altogether orthodox; but alas, what a cold, joyless, forsaken Christianity! A Christianity without the friend, the personal friend, Jesus!

For such he is to the one who flees to him in childlike trustful simplicity. If he is the Son of God, then he is also Mary’s Son, a human being like you and me. He wept over Jerusalem, and with unspeakable sorrow he both saw and foretold the dreadful judgment that would come upon David’s city and the chosen people, indeed upon the whole world, for the sake of obstinacy and unbelief. But the compassion of his heart was not confined to these great events of world history. He had care for the distress and need of each individual, in small things as in great. He had deep compassion on the people who had been with him for several days without getting food; he took thought that they should not perish on the way, and he fed them. He stopped when he saw the weeping widow in Nain, had deep compassion on her, and gave her back her son. And when he saw Mary weeping, and those who were with her weeping over the death of Lazarus, then it is written of him that he was deeply moved in spirit, indeed shaken, and that he also wept.

Jesus wept! Such a friend is he; you need not be ashamed or afraid to go to him with everything that lies upon your heart; for he both will and can feel with you; he can sing with the glad and weep with the sorrowing. Therefore it is written of him in Hebrews 2:16-18, that he does not come to the aid of angels, but of the offspring of Abraham, and that he had to be made like his brothers in all things, and since he has suffered and has himself been tempted, he can come to the help of those who are tempted.

But if he can so truly and deeply feel with us and understand us in all the movements of our heart, then he is also ready to give up all that is his in order to help us. That was what he did for Martha and Mary. "No one has greater love than this, that one lay down his life for his friends," he himself says; this rule he has followed for us all on the Cross, when we were enemies. How much more readily, then, does he not give all for his friends!

"Where have you laid him?" he asked; by then he had already taken counsel with his Father in Heaven and formed his decision as to what he would do. He knew that the raising of Lazarus would be like the last drop in the cup of bitterness and offense which his words and deeds had poured out for the leaders of the Jews. They would resolve upon his death, and the high priest Caiaphas would even adorn their evil deed with the hypocritical appearance of godliness. But nothing could hold Jesus back from helping his friends in their sorrow and distress. His heart was full; and as he followed them to the tomb, he wept, not only over sin and its painful consequences, death and sorrow, but also out of true and deep compassion for the mourners, whom he therefore also willed to help.

Then, as now, there were many kinds of astonishment at this, which pride and wisdom often call weakness, that Jesus wept. "See how he loved him," said some. It no doubt seemed strange to them that the great prophet, who had aroused attention and in part even stir throughout the whole land, could concern himself with such small things and in such a degree let himself be carried away by his personal feelings.

"Could not he who opened the blind man’s eyes also have caused this man not to die?" asked others.

These were the same who, when they had brought him to the Cross, cried out, "He helped others; himself he cannot help; let him come down from the Cross, and we will believe."

That is the speech of unbelief and heartlessness, which cannot understand that when a man has power, he should not use it for his own advantage. Love is to them a hidden and unknown thing. That Jesus had power to heal diseases and to raise the dead, and yet should not use it to spare himself sorrow and suffering, was to them wholly incomprehensible. That God’s own Son became man and died the death of scorn for sinners is to unbelief, to this very day, incomprehensible and offensive.

But his power, then as now, Jesus used to relieve and heal the sorrow of those who turned to him in trust. He did far more than they dared ask or understand. Such a friend is he. "Lord, if you had been here before, my brother would not have died," and, "he already stinks; he has now lain four days in the tomb," thus the two sisters spoke to their friend. His love and power were greater than both their trust and their fear. Only believe, Mary; do not fear, Martha! and you shall see the glory of God. Thus he went to the tomb and cried, "Lazarus, come out!" and he came out. The sisters received their brother again, and the wound of their hearts was healed. The love of Jesus had done it all. Such a friend is he.

But he did not do it apart from the will of the Father. His love was, like all true love, divine. It has its source and its rule in God. This same Jesus, who says of himself, "All power in Heaven and on Earth is given to me," goes as an obedient Son to the Father, asks him for everything, and does nothing apart from his will. "Father, I thank you that you have heard me," he cried, before he stepped to the tomb and raised the dead. Power and pride most often go together among men. Omnipotence and humility go side by side in the Son of Man and form the living foundation of his love. Such a friend is Jesus.

Are you seeking such a friend? Oh, here he is at the door of your heart; he knocks and asks whether he may come in and be for you what he was for Martha and Mary. Have you found him, oh, then hold fast to him in life and death, and fear nothing.

When sorrow and sickness and distress, all the painful consequences of sin, go like a flood over your soul, you have a fully faithful friend; keep near to him, speak with him childlike and confidently, complain to him of all your distress, count nothing too small to lay before him; he both can and will help far, far beyond what you might even dare to ask or understand; come like Martha with all your thoughts and all your doubts, come like Mary with all the singleness of your heart and your trustful spirit, he shall relieve your sorrow and heal your soul and, in the resurrection of the dead, give you an eternal and unfading joy and blessedness.

Mary, Martha, do you believe this?

"I know a friend who smiles on me
When I am set in sorrow and distress;
He fails me not, in life and death,
In the storm I rest upon his breast.
A friend so mighty and so fair,
The fairest of Adam’s race,
A friend who never fails."




%FIXME: https://hymnary.org/text/jeg_veed_en_vei_saa_fuld_af_traengsel

Meter: 9.8.8.9.8.8.7

%Candidate tune: https://hymnary.org/tune/at_peace_with_god_how_great_the_blessing



1 Jeg veed en Vei, saa fuld af Trængsel,
En Taarevei saa tung og trang;
Men henad Veien lyder Sang
Om Troens Kamp og Himlens Længsel;
Det er den Vei, de Kristne gaa
For Livsens Krone at opnaa
Igjennem Kamp og Trængsel.

2 Jeg veed en Ven, som til mig smiler,
Naar jeg er stedt i Sorg og Nød;
Han sviger ei i Liv og Død;
I Stormen ved hans Bryst jeg hviler;
En Ven saa mægtig og saa skjøn,
Den deiligste af Adams Kjøn,
En Ven, som aldrig sviger.

3 Jeg veed en Dragt, mod hvilken blegner
Al Verdens Pragt og Herlighed,
Som hver en synder smykkes med,
Naar Gud ham sønnens Blod tilregner;
En Bryllupsklædning ren og prud,
Som skjuler Synderne for Gud,
En Dragt for Brudeskaren.

% “set Fængsel” should be “sit Fængsel.”

4 Jeg veed en Stund, naar bud der kommer
Om Hvile efter Dagens Strid,
Da byttes Verdens Vintertid
Med Evighedens glade Sommer;
Da foldes Blomsterknoppen ud,
Og Rosen blomstrer sjønt hos Gud,
Da har den sprængt set Fængsel.

5 Jeg veed et Hjem, hvor Børneskaren
Fra alle Kanter samles skal,
Det er den lyse Himmelsal,
Til hvilken Kristus er opfaren,
Der samles de fra Øst og Vest,
fra Nord og Syd til Bryllupsfest,
Og Glæden faar ei Ende.


HYMN TRANSLATION

I know a way so full of tribulation, 
A way of tears so heavy and so narrow; But along that way there sounds a song Of faith’s good fight and Heaven’s longing; It is the way the Christians go To win the crown of life Through conflict and tribulation.

I know a friend who smiles upon me
When I am set in sorrow and distress; He does not fail, in life and death; In the storm I rest upon his breast. A friend so mighty and so fair, The fairest of Adam’s race, A friend who never fails.

I know a garment, before which pales All the world’s splendor and glory, With which every sinner is adorned When God reckons to him the Son’s blood; A wedding-garment, pure and comely, Which hides sins before God, A garment for the bridal company.

I know an hour when word shall come 
Of rest after the day’s strife, When the world’s wintertime is exchanged For the glad summer of eternity; Then the flower-bud unfolds, And the rose blooms fair with God, Then it has burst its prison.

I know a home where the band of children 
Shall gather from every quarter; It is the bright hall of Heaven, To which Christ has ascended; There they gather from east and west, From north and south, to the marriage-feast, And joy shall have no end.